\documentclass[12pt,a4paper]{article}
\usepackage{../latex/dsdlab-report}

\begin{document}
	\reporttitlepage{۳}{توصیف جریان داده}

	\section*{۳.۱ مقدمه و اهداف}
	در این آزمایش، توصیف جریان داده در زبان توصیف سخت‌افزار \lr{Verilog} با دستور پیوستهٔ \lr{assign} بررسی می‌شود. بخش نخست شامل طراحی سلسله‌مراتبی مقایسه‌کنندهٔ آبشاری چهاربیتی و بخش دوم شامل طراحی یک مقایسه‌کنندهٔ پی‌درپی دارای ورودی بازنشانی است.

	\subsection*{روند انجام آزمایش}
	برای هر بخش، ابتدا پروژه‌ای جداگانه ساخته و پرونده‌های توصیف سخت‌افزار به آن افزوده شدند. سپس پیمانهٔ سطح بالا تعیین، خانوادهٔ قطعه انتخاب و طرح همگردانی و سنتز شد. در بخش آبشاری، پیمانهٔ یک‌بیتی پیش از پیمانهٔ چهاربیتی ساخته شد؛ در بخش پی‌درپی، همهٔ منطق در یک پیمانه قرار گرفت. پس از ساخت و اجرای میزهای آزمون در \lr{ModelSim}، خروجی‌ها با نتیجهٔ مرجع مقایسه شدند و در پایان، نمای سطح انتقال ثبات‌ها و گزارش مصرف منابع هر دو طرح بررسی شد.

	\section*{۳.۲ طراحی و کد توصیف سخت‌افزار}
	\subsection*{الف) مقایسه‌کنندهٔ آبشاری چهاربیتی}
	واحد پایه، یک مقایسه‌کنندهٔ یک‌بیتی آبشاری با سه ورودی وضعیّت از مرحلهٔ پرارزش‌تر است. اگر مراحل قبلی برابری را گزارش کرده باشند، بیت جاری می‌تواند نتیجه را به بزرگ‌تر یا کوچک‌تر تغییر دهد؛ در غیر این صورت، نتیجهٔ قطعی مرحلهٔ قبلی بدون تغییر عبور می‌کند. روابط داده‌ای این مرحله عبارت‌اند از:
	\[
	\begin{aligned}
	G_{\mathrm{next}}&=G_{\mathrm{prev}}\lor(E_{\mathrm{prev}}\land A\land\overline{B}),\\
	E_{\mathrm{next}}&=E_{\mathrm{prev}}\land\overline{A\oplus B},\\
	L_{\mathrm{next}}&=L_{\mathrm{prev}}\lor(E_{\mathrm{prev}}\land\overline{A}\land B).
	\end{aligned}
	\]
	هر سه رابطه تنها با دستورهای پیوستهٔ \lr{assign} توصیف شده‌اند.

	\shot{screenshots/Cascadable_Comparator/module_compare_stage_code.png}{کد جریان داده‌ی مرحله‌ی یک‌بیتی مقایسه‌کننده}

	برای ساخت مقایسه‌کنندهٔ چهاربیتی، چهار نمونه از این مرحله از بیت ۳ تا بیت ۰ به‌صورت آبشاری متّصل شدند. وضعیّت اوّلیّهٔ مرحلهٔ پرارزش \(G=0,E=1,L=0\) است و خروجی هر مرحله به ورودی وضعیّت مرحلهٔ بعد می‌رود. در نتیجه، نخستین اختلاف از سمت بیت پرارزش نتیجهٔ نهایی را تعیین می‌کند.

	\shot{screenshots/Cascadable_Comparator/topfour_bit_comparator_code_1.png}{تعریف ورودی‌ها، خروجی‌ها و مرحله‌ی بیت ۳}
	\shot{screenshots/Cascadable_Comparator/topfour_bit_comparator_code_2.png}{اتّصال مرحلهٔ بیت ۳ به مرحلهٔ بیت ۲}
	\shot{screenshots/Cascadable_Comparator/topfour_bit_comparator_code_3.png}{اتّصال مراحل بیت ۱ و صفر و تولید خروجی نهایی}

	نمای سطح انتقال ثبات‌ها دقیقاً چهار نمونه از پیمانهٔ \lr{compare\_stage} را به‌صورت متوالی نشان می‌دهد و بنابراین ساختار سلسله‌مراتبی موردنظر دستورکار حفظ شده است.

	\shot{screenshots/Cascadable_Comparator/RTL_photo.png}{نمای سطح انتقال ثبات‌های مقایسه‌کنندهٔ آبشاری چهاربیتی}

	\subsection*{ب) مقایسه‌کنندهٔ پی‌درپی}
	مقایسه‌کنندهٔ پی‌درپی پس از فعّال‌شدن ورودی بازنشانی \lr{reset} در حالت برابر، یعنی \(G=0,E=1,L=0\)، قرار می‌گیرد. بیت‌های دو عدد از پرارزش‌ترین بیت به کم‌ارزش‌ترین بیت وارد می‌شوند. تا وقتی \(E=1\) است، هر بیت جدید می‌تواند رابطه را تعیین کند؛ پس از مشاهدهٔ نخستین اختلاف، وضعیّت بزرگ‌تر یا کوچک‌تر حفظ می‌شود. روابط حالت بعد همان روابط مرحلهٔ آبشاری‌اند، امّا نتیجه باید بین پالس‌های ساعت نگهداری شود.

	برای رعایت محدودیّت استفادهٔ انحصاری از \lr{assign}، نگهداری وضعیّت با شش شبکهٔ بازخوردی در دو مرحلهٔ اصلی و تابع انجام شده است. در یک سطح ساعت، حالت بعد در سیگنال‌های مرحلهٔ اصلی با پیشوند \lr{master\_} ذخیره می‌شود و در سطح دیگر به خروجی منتقل می‌گردد. بازنشانی نیز به‌صورت غیرهم‌زمان مقادیر اوّلیّه را تحمیل می‌کند. مقایسه‌کنندهٔ پی‌درپی در یک پیمانه و بدون نمونه‌سازی هیچ پیمانهٔ دیگری پیاده‌سازی شده است؛ بنابراین تمام منطق آن، مطابق محدودیّت دستورکار، بدون توصیف رفتاری و فقط با توصیف جریان داده ساخته شده است.

	\shot{screenshots/Serial_Comparator/serail_bit_comparator_code_1.png}{تعریف رابط و روابط حالت بعد مقایسه‌کنندهٔ پی‌درپی}
	\shot{screenshots/Serial_Comparator/serail_bit_comparator_code_2.png}{پیاده‌سازی نگهداری حالت و بازنشانی تنها با \lr{assign}}
	\shot{screenshots/Serial_Comparator/RTL_view.png}{نمای سطح انتقال ثبات‌ها و عناصر حافظهٔ استنتاج‌شده برای مقایسه‌کنندهٔ پی‌درپی}

	\section*{۳.۳ میز آزمون و شبیه‌سازی}
	\subsection*{الف) آزمون مقایسه‌کننده‌ی آبشاری}
	میز آزمون آبشاری تمام \(16\times16=256\) ترکیب ممکن دو ورودی چهاربیتی را پیمایش می‌کند. برای هر بردار، سه خروجی مدار با عملگرهای مرجع \lr{>}، \lr{==} و \lr{<} مقایسه شده و هر مغایرت در شمارندهٔ خطای \lr{error\_count} ثبت می‌شود.

	\shot{screenshots/Cascadable_Comparator/test_bench_code_1.png}{تعریف میز آزمون و نمونه‌سازی مقایسه‌کنندهٔ چهاربیتی}
	\shot{screenshots/Cascadable_Comparator/test_bench_code_2.png}{حلقه‌های آزمون فراگیر و مقایسه با نتیجه‌ی مرجع}
	\shot{screenshots/Cascadable_Comparator/test_bench_code_3.png}{گزارش قبولی یا خطای آزمون فراگیر}

	اجرای میز آزمون در \lr{ModelSim} با همگردانی سه پیمانه و افزودن سیگنال‌ها به پنجرهٔ موج انجام شد. خروجی پایانه، قبولی تمام ۲۵۶ حالت و صفر بودن شمارندهٔ خطا را نشان می‌دهد. شکل موج کامل و نمای پالایش‌شده نیز ویژگی تک‌فعّال خروجی‌ها را تأیید می‌کنند؛ در هر لحظه دقیقاً یکی از سه رابطهٔ بزرگ‌تر، برابر یا کوچک‌تر فعّال است.

	\shot{screenshots/Cascadable_Comparator/how_to_run_in_modelsim.png}{فرمان‌های اجرای میز آزمون در \lr{ModelSim}}
	\shot{screenshots/Cascadable_Comparator/Cascadable_terminal_modelsim.png}{پیام قبولی همه‌ی ۲۵۶ حالت آزمون}
	\shot{screenshots/Cascadable_Comparator/Cascadable_wave_all_signals.png}{شکل موج همه‌ی سیگنال‌های مقایسه‌کننده‌ی آبشاری}
	\shot{screenshots/Cascadable_Comparator/Cascadable_wave_only_necessary_signals.png}{شکل موج ورودی‌ها و خروجی‌های اصلی}
	\shot{screenshots/Cascadable_Comparator/Cascadable_wave_only_necessary_signals_zoom.png}{نمای بزرگ‌شده‌ی تغییر خروجی‌ها برای چند بردار نخست}

	\subsection*{ب) آزمون مقایسه‌کنندهٔ پی‌درپی}
	میز آزمون پی‌درپی یک ساعت با نیم‌دورهٔ ۵ نانوثانیه تولید می‌کند. وظیفهٔ \lr{reset\_comparator} درستی حالت اوّلیّه را کنترل می‌کند و وظیفهٔ \lr{compare\_words} بیت‌های دو کلمهٔ چهاربیتی را از بیت ۳ تا صفر، در لبهٔ منفی، اعمال می‌کند. خروجی پس از لبهٔ مثبت با یک الگوی مرجع مقایسه می‌شود؛ الگوی مرجع نیز فقط تا زمان برابر بودن پیشوند، اجازهٔ تغییر نتیجه را می‌دهد.

	\shot{screenshots/Serial_Comparator/test_bench_code_1.png}{تعریف سیگنال‌ها، مدار تحت آزمون و مولّد ساعت}
	\shot{screenshots/Serial_Comparator/test_bench_code_2.png}{وظیفهٔ بازنشانی و کنترل حالت اوّلیّه}
	\shot{screenshots/Serial_Comparator/test_bench_code_3.png}{تعریف وظیفهٔ مقایسه و ارسال بیت‌ها از بیت پرارزش}
	\shot{screenshots/Serial_Comparator/test_bench_code_4.png}{به‌روزرسانی الگوی مرجع و مقایسهٔ خروجی‌ها}
	\shot{screenshots/Serial_Comparator/test_bench_code_5.png}{گزارش نتیجهٔ هر بیت و پایان وظیفه}
	\shot{screenshots/Serial_Comparator/test_bench_code_6.png}{بردارهای آزمون برابر، بزرگ‌تر و کوچک‌تر}
	\shot{screenshots/Serial_Comparator/test_bench_code_7.png}{گزارش نهایی قبولی یا تعداد خطاها}

	شش زوج آزمون شامل حالت برابر و نمونه‌های بزرگ‌تر و کوچک‌تر اعمال شدند. متن پایانه نتیجهٔ هر بیت را نشان می‌دهد و در پایان، پیام قبولی همهٔ آزمون‌ها چاپ شده است. شکل موج‌ها نیز بازنشانی به \(010\)، دریافت پی‌درپی بیت‌ها و تثبیت نتیجه پس از نخستین اختلاف معنادار را نشان می‌دهند.

	\shot{screenshots/Serial_Comparator/terminal_passed_all_tests.png}{خروجی پایانه و قبولی تمام آزمون‌های پی‌درپی}
	\shot{screenshots/Serial_Comparator/test_all_waves.png}{شکل موج کامل سیگنال‌های داخلی و میز آزمون}
	\shot{screenshots/Serial_Comparator/test_necessary_waves.png}{شکل موج سیگنال‌های ضروری مقایسه‌کنندهٔ پی‌درپی}

	\section*{۳.۴ سنتز و پیاده‌سازی نرم‌افزاری}
	هر دو طرح برای قطعهٔ \lr{EP2C5T144C6} از خانوادهٔ \lr{Cyclone II} سنتز شدند. مقایسه‌کنندهٔ آبشاری ۹ تابع ترکیبی و ۱۱ پایه، و مقایسه‌کنندهٔ پی‌درپی ۹ تابع ترکیبی و ۷ پایه مصرف می‌کند؛ در هر دو مورد، مصرف عناصر منطقی کمتر از یک درصد ظرفیّت تراشه است.

	\shot{screenshots/Cascadable_Comparator/flow_summary.png}{خلاصه‌ی منابع مقایسه‌کننده‌ی آبشاری}
	\shot{screenshots/Serial_Comparator/flow_summary.png}{خلاصهٔ منابع مقایسه‌کنندهٔ پی‌درپی}

	همگردانی هر دو پروژه با صفر خطا تکمیل شد. طرح آبشاری به‌دلیل ترکیبی بودن ساعت ندارد و هشدارهای تحلیل زمان برای آن طبیعی است. برای طرح پی‌درپی قید زمانی مستقلی تعریف نشده است؛ ازاین‌رو، هشدارهای تحلیل زمان و کمینهٔ پهنای پالس باقی مانده‌اند. همچنین در پرونده‌های پیکربندی پروژه تخصیص پایه‌ای برای کلیدها، ساعت و چراغ‌های برد ثبت نشده است. با توجّه به برگزاری مجازی آزمایشگاه، شبیه‌سازی و سنتز نرم‌افزاری مبنای ارزیابی این آزمایش بوده و ادّعایی دربارهٔ اجرای واقعی روی برد مطرح نمی‌شود.

	\shot{screenshots/Cascadable_Comparator/successful_compile.png}{همگردانی موفّق مقایسه‌کنندهٔ آبشاری}
	\shot{screenshots/Serial_Comparator/successful_compile.png}{همگردانی موفّق مقایسه‌کنندهٔ پی‌درپی}

	\section*{۳.۵ نتایج}
	\begin{itemize}
		\item مقایسه‌کنندهٔ چهاربیتی با اتّصال سلسله‌مراتبی چهار مرحلهٔ یک‌بیتی و صرفاً با توصیف جریان داده ساخته شد.
		\item آزمون فراگیر تمام ۲۵۶ زوج ورودی بدون خطا پایان یافت.
		\item مقایسه‌کنندهٔ پی‌درپی نتیجهٔ مقایسهٔ پیشوند دریافت‌شده را در هر پالس نگه داشت و پس از نخستین اختلاف، وضعیّت نهایی را حفظ کرد.
		\item حالت بازنشانی \(G=0,E=1,L=0\) و هر شش زوج آزمون پی‌درپی با موفّقیّت بررسی شدند.
		\item هر دو پروژه بدون خطای همگردانی و با مصرف کمتر از یک درصد منابع منطقی سنتز شدند.
	\end{itemize}

	\section*{نتیجه‌گیری}
	این آزمایش نشان داد که توصیف جریان داده هم برای مدار ترکیبی سلسله‌مراتبی و هم، با استفاده از بازخورد کنترل‌شده، برای مدار ترتیبی قابل‌استفاده است. زنجیرهٔ مقایسه از بیت پرارزش به کم‌ارزش، منطق اولویت عددی را مستقیماً پیاده‌سازی کرد و نسخهٔ پی‌درپی همان منطق را در طول زمان نگه داشت. میز آزمون فراگیر بخش اوّل، بررسی بیت‌به‌بیت بخش دوم و سنتز موفّق هر دو طرح، درستی عملکرد را مطابق اهداف آزمایش مجازی تأیید کردند.
\end{document}
